\documentclass[12pt]{article}
\usepackage[letterpaper, margin=1.5cm]{geometry}
\usepackage{amsmath}
\usepackage{amssymb}
\setlength\parindent{0pt}

\begin{document}
Se define:
\begin{align} 
    \pmb{\dot\varepsilon}^p &= \gamma \frac{\partial f\left(\pmb{\sigma}, \pmb{q}\right)}{\partial\pmb{\sigma}} \\
    \pmb{\dot \alpha} &= \gamma \frac{\partial f\left(\pmb{\sigma}, \pmb{q}\right)}{\partial\pmb{q}} 
    % \pmb{\dot q} &= -\gamma \mathbb{D} \frac{\partial f\left(\pmb{\sigma}, \pmb{q}\right)}{\partial\pmb{q}}
\end{align}
Con lo anterior se evalúa la derivada temporal de los esfuerzos $\pmb{q}$ producidos por las variables internas $\pmb{\alpha}$.
\[ \pmb{\dot q} = -\nabla \pmb{q}\left(\pmb{\alpha}\right) \pmb{\dot\alpha} = -\gamma \mathbb{D} \frac{\partial f\left(\pmb{\sigma}, \pmb{q}\right)}{\partial\pmb{q}}\]

Discretizando por euler implicito
\begin{align} 
    \pmb{\varepsilon}_{n+1}^p &= \pmb{\varepsilon}_{n}^p + \Delta\gamma\partial_{\pmb{\sigma}} f\left(\pmb{\sigma}_{n+1}, \pmb{q}_{n+1}\right) \\
    % \pmb{q}_{n+1} &= \pmb{q}_{n} - \Delta\gamma\mathbb{D}\partial_{\pmb{q}} f\left(\pmb{\sigma}_{n+1}, \pmb{q}_{n+1}\right)
    \pmb{\alpha}_{n+1} &= \pmb{\alpha}_{n} + \Delta\gamma\partial_{\pmb{q}} f\left(\pmb{\sigma}_{n+1}, \pmb{q}_{n+1}\right)
\end{align}

De esta forma se obtiene el sistema de ecuaciones siguiente:
\begin{align} 
\pmb{\varepsilon}_{n+1} &= \pmb{\varepsilon}_n + \Delta\pmb{\varepsilon} \\
\pmb{\sigma}_{n+1} &= \mathbb{C} : \left( \pmb{\varepsilon}_{n+1} - \pmb{\varepsilon}_{n+1}^p \right) \\
\pmb{\varepsilon}_{n+1}^p &= \pmb{\varepsilon}_{n}^p + \Delta\gamma\partial_{\pmb{\sigma}} f\left(\pmb{\sigma}_{n+1}, \pmb{q}_{n+1}\right)\\
\pmb{\alpha}_{n+1} &= \pmb{\alpha}_{n} + \Delta\gamma\partial_{\pmb{q}} f\left(\pmb{\sigma}_{n+1}, \pmb{q}_{n+1}\right)
\end{align}

Se definen las componentes de prueba
\begin{align} 
    \pmb{\varepsilon}_{n+1}^{e\textnormal{ trial}} &= \pmb{\varepsilon}_{n+1} - \pmb{\varepsilon}_{n}^p \\
    \pmb{\sigma}_{n+1}^{\textnormal{trial}} &= \mathbb{C} : \pmb{\varepsilon}_{n+1}^{e\textnormal{ trial}} \\
    \pmb{q}_{n+1}^{\textnormal{trial}} &= \pmb{q}_{n+1} \\
    f_{n+1}^{\textnormal{trial}} &= f\left(\pmb{\sigma}_{n+1}^{\textnormal{trial}}, \pmb{q}_{n+1}^{\textnormal{trial}}\right)
\end{align}

Para simplificar las ecuaciones se define la función de fluencia y las demás variables:
\begin{align}
    f\left(\pmb{\sigma}_{n+1}, \bar\varepsilon^p_{n+1}\right) &= \sigma_{eq}\left(\pmb{\sigma}_{n+1}\right) - Y\left(\bar\varepsilon^p_{n+1}\right) \label{eqn:yield}\\
    \pmb{\varepsilon}_{n+1}^p &= \pmb{\varepsilon}^p_{n} + \Delta\gamma\partial_{\pmb{\sigma}}\sigma_{eq}\left(\pmb{\sigma}_{n+1}\right)  \label{eqn:plas_strain_imp}\\
    \pmb{\alpha}_{n+1} &= \bar\varepsilon^p_{n+1} = \bar\varepsilon^p_{n} + \Delta\gamma  \label{eqn:eff_plas_strain_imp}\\
    \pmb{q}_{n+1} &= Y(\bar\varepsilon^p_{n+1}) \\
\end{align}

En base a lo anterior se definen los dos residuos para poder hacer Newton-Rhapson. El primero $\pmb{R}_{n+1}^{\varepsilon^p}$ respecto de la deformación plástica (ecuación~\ref{eqn:plas_strain_imp}), el segundo $\pmb{R}_{n+1}^{f}$ respecto de la función de fluencia (ecuación~\ref{eqn:yield}):
\begin{align}
    \pmb{R}_{n+1}^{\varepsilon^p} &= -\pmb{\varepsilon}_{n+1}^p + \pmb{\varepsilon}_{n}^p + \Delta\gamma\partial_{\pmb{\sigma}}\sigma_{eq}\left(\pmb{\sigma}_{n+1}\right) \\
    \pmb{R}_{n+1}^{f} &= \sigma_{eq}\left(\pmb{\sigma}_{n+1}\right) - Y\left(\bar\varepsilon^p_{n+1}\right)
\end{align}

Como $\pmb{\varepsilon}^p_{n+1}$ es constante en el algoritmo de retorno plástico, se puede obtener la siguiente expresión en las iteraciones de Newton-Rhapson:
\[
\left(\pmb{\sigma}_{n+1}^{(k)} + \Delta\pmb{\sigma}_{n+1}^{(k)}\right) = \mathbb{C}:\left[\pmb{\varepsilon}_{n+1} - \left(\pmb{\varepsilon}_{n+1}^{p^{(k)}} + \Delta\pmb{\varepsilon}_{n+1}^{p^{(k)}}\right)\right] \implies \Delta\pmb{\sigma}_{n+1}^{(k)} = -\mathbb{C}:\Delta\pmb{\varepsilon}_{n+1}^{p^{(k)}}
\]\[
    \Delta\pmb{\varepsilon}_{n+1}^{p^{(k)}} = -\mathbb{C}^{-1}:\Delta\pmb{\sigma}_{n+1}^{(k)}
\]

Ahora se linealiza el sistema por medio de los residuos
\[\renewcommand\arraystretch{1.6}
    \begin{bmatrix}
        \frac{\partial\pmb{R}_{n+1}^{\varepsilon^p}}{\partial \pmb{\varepsilon}^p_{n+1}} & \frac{\partial\pmb{R}_{n+1}^{\varepsilon^p}}{\partial \Delta\gamma_{n+1}} \\ 
        \frac{\partial\pmb{R}_{n+1}^{f}}{\partial \pmb{\varepsilon}^p_{n+1}} & \frac{\partial\pmb{R}_{n+1}^{f}}{\partial \Delta\gamma_{n+1}} \\
    \end{bmatrix}
    \begin{bmatrix}
        \Delta\pmb{\varepsilon}^p_{n+1} \\
        \Delta^2\gamma_{n+1} \\
    \end{bmatrix} = -
    \begin{bmatrix}
        \pmb{R}_{n+1}^{\varepsilon^p} \\
        \pmb{R}_{n+1}^{f} \\
    \end{bmatrix} 
\]
Tomando $\pmb{\Xi} = \left[\mathbb{C}^{-1} + \Delta\gamma\partial^2_{\sigma\sigma}\sigma_{eq}\left(\pmb{\sigma}_{n+1}\right)\right]^{-1}$ 
\[\renewcommand\arraystretch{1.6}
    \begin{bmatrix}
        \pmb{\Xi}^{-1}:\Delta\pmb{\sigma}_{n+1}:\left[\Delta\pmb{\varepsilon}^p_{n+1}\right]^{-1} & \partial_{\sigma}\sigma_{eq}\left(\pmb{\sigma}_{n+1}\right) \\ 
        \partial_{\sigma}\sigma_{eq}\left(\pmb{\sigma}_{n+1}\right):\Delta\pmb{\sigma}_{n+1}:\left[\Delta\pmb{\varepsilon}^p_{n+1}\right]^{-1} & -\partial_{\Delta\gamma}Y\left(\bar\varepsilon^p_{n+1}\right)
    \end{bmatrix}
    \begin{bmatrix}
        \Delta\pmb{\varepsilon}^p_{n+1} \\
        \Delta^2\gamma_{n+1} \\
    \end{bmatrix} = -
    \begin{bmatrix}
        \pmb{R}_{n+1}^{\varepsilon^p} \\
        \pmb{R}_{n+1}^{f} \\
    \end{bmatrix}
\]
De la primera ecuación se despeja $\Delta\pmb{\sigma}_{n+1}$:
\[ \pmb{\Xi}^{-1}:\Delta\pmb{\sigma}_{n+1} + \partial_{\sigma}\sigma_{eq}\left(\pmb{\sigma}_{n+1}\right)\Delta^2\gamma + \pmb{R}_{n+1}^{\varepsilon^p} = 0\]
\[ \Delta\pmb{\sigma}_{n+1} = \pmb{\Xi}:\left[-\partial_{\sigma}\sigma_{eq}\left(\pmb{\sigma}_{n+1}\right)\Delta^2\gamma - \pmb{R}_{n+1}^{\varepsilon^p} \right]\]

De la segunda ecuación
\[ \partial_{\sigma}\sigma_{eq}\left(\pmb{\sigma}_{n+1}\right):\Delta\pmb{\sigma}_{n+1} - \partial_{\Delta\gamma}Y\left(\bar\varepsilon^p_{n+1}\right)\Delta^2\gamma_{n+1} + \pmb{R}_{n+1}^{f} = 0\]
\[ -\partial_{\sigma}\sigma_{eq}\left(\pmb{\sigma}_{n+1}\right):\pmb{\Xi}:\partial_{\sigma}\sigma_{eq}\left(\pmb{\sigma}_{n+1}\right)\Delta^2\gamma
   -\partial_{\sigma}\sigma_{eq}\left(\pmb{\sigma}_{n+1}\right):\pmb{\Xi}:\pmb{R}_{n+1}^{\varepsilon^p} 
   - \partial_{\Delta\gamma}Y\left(\bar\varepsilon^p_{n+1}\right)\Delta^2\gamma_{n+1} + \pmb{R}_{n+1}^{f} = 0\]

Factorizando por $\Delta^2\gamma$ y despejando
\[
\Delta^2\gamma = \frac{\pmb{R}_{n+1}^{f} - \partial_{\sigma}\sigma_{eq}\left(\pmb{\sigma}_{n+1}\right):\pmb{\Xi}:\pmb{R}_{n+1}^{\varepsilon^p}}
{\partial_{\sigma}\sigma_{eq}\left(\pmb{\sigma}_{n+1}\right):\pmb{\Xi}:\partial_{\sigma}\sigma_{eq}\left(\pmb{\sigma}_{n+1}\right) + \partial_{\Delta\gamma}Y\left(\bar\varepsilon^p_{n+1}\right)}
\]


\section*{Tangente consistente}
\[\pmb{\sigma}_{n+1} = \mathbb{C}:\left(\pmb{\varepsilon}_{n+1} - \pmb{\varepsilon}_{n+1}^p\right) \]
\begin{equation} \label{eqn:hooke_differential}
    d\pmb{\sigma}_{n+1} = \mathbb{C}:\left(d\pmb{\varepsilon}_{n+1} - d\pmb{\varepsilon}_{n+1}^p\right)
\end{equation}
Luego se diferencia el residuo de deformación plástico $\pmb{R}_{n+1}^{\varepsilon^p}$
\[\pmb{R}_{n+1}^{\varepsilon^p} = 0 = -\pmb{\varepsilon}_{n+1}^p + \pmb{\varepsilon}_{n}^p + \Delta\gamma\partial_{\pmb{\sigma}}\sigma_{eq}\left(\pmb{\sigma}_{n+1}\right) \]
\[d\pmb{\varepsilon}_{n+1}^p = d\Delta\gamma\partial_{\pmb{\sigma}}\sigma_{eq}\left(\pmb{\sigma}_{n+1}\right) + \Delta\gamma\partial_{\pmb{\sigma}\pmb{\sigma}}^2 \sigma_{eq}\left(\pmb{\sigma}_{n+1}\right) : d\pmb{\sigma}_{n+1} \]

El término $d\pmb{\varepsilon}_{n+1}^p$ se reemplaza en la ecuación~\ref{eqn:hooke_differential}. 
\[d\pmb{\sigma}_{n+1} = \mathbb{C}:\left(d\pmb{\varepsilon}_{n+1} - d\Delta\gamma\partial_{\pmb{\sigma}}\sigma_{eq}\left(\pmb{\sigma}_{n+1}\right) - \Delta\gamma\partial_{\pmb{\sigma}\pmb{\sigma}}^2 \sigma_{eq}\left(\pmb{\sigma}_{n+1}\right) : d\pmb{\sigma}_{n+1}\right)\]
\[\mathbb{C}^{-1}:d\pmb{\sigma}_{n+1} = d\pmb{\varepsilon}_{n+1} - d\Delta\gamma\partial_{\pmb{\sigma}}\sigma_{eq}\left(\pmb{\sigma}_{n+1}\right) - \Delta\gamma\partial_{\pmb{\sigma}\pmb{\sigma}}^2 \sigma_{eq}\left(\pmb{\sigma}_{n+1}\right) : d\pmb{\sigma}_{n+1}\]
\[\pmb{\Xi}^{-1}:d\pmb{\sigma}_{n+1} = d\pmb{\varepsilon}_{n+1} - d\Delta\gamma\partial_{\pmb{\sigma}}\sigma_{eq}\left(\pmb{\sigma}_{n+1}\right)\]
\[d\pmb{\sigma}_{n+1} = \pmb{\Xi}:\left[d\pmb{\varepsilon}_{n+1} - d\Delta\gamma\partial_{\pmb{\sigma}}\sigma_{eq}\left(\pmb{\sigma}_{n+1}\right)\right]\]

Luego se utiliza la función de fluencia y se diferencia.
\[ f(\pmb{\sigma}_{n+1}, \Delta\gamma_{n+1}) = 0 = \sigma_{eq}\left(\pmb{\sigma}_{n+1}\right) - Y\left(\bar\varepsilon^p_{n+1}\right)\]
\[ 0 = \partial_{\pmb{\sigma}}\sigma_{eq}\left(\pmb{\sigma}_{n+1}\right):d\pmb{\sigma}_{n+1} - \partial_{\Delta\gamma} Y\left(\bar\varepsilon^p_{n+1}\right) d\Delta\gamma_{n+1}\]
\[ 0 = \partial_{\pmb{\sigma}}\sigma_{eq}\left(\pmb{\sigma}_{n+1}\right):\pmb{\Xi}:\left[d\pmb{\varepsilon}_{n+1} - d\Delta\gamma\partial_{\pmb{\sigma}}\sigma_{eq}\left(\pmb{\sigma}_{n+1}\right)\right] - \partial_{\Delta\gamma} Y\left(\bar\varepsilon^p_{n+1}\right) d\Delta\gamma_{n+1}\]
\[ 0 = \partial_{\pmb{\sigma}}\sigma_{eq}\left(\pmb{\sigma}_{n+1}\right):\pmb{\Xi}:d\pmb{\varepsilon}_{n+1} - d\Delta\gamma\left[\partial_{\pmb{\sigma}}\sigma_{eq}\left(\pmb{\sigma}_{n+1}\right):\pmb{\Xi}:\partial_{\pmb{\sigma}}\sigma_{eq}\left(\pmb{\sigma}_{n+1}\right) + \partial_{\Delta\gamma} Y\left(\bar\varepsilon^p_{n+1}\right)\right]\]
\[ d\Delta\gamma = \frac{\partial_{\pmb{\sigma}}\sigma_{eq}\left(\pmb{\sigma}_{n+1}\right):\pmb{\Xi}:d\pmb{\varepsilon}_{n+1} }{\partial_{\pmb{\sigma}}\sigma_{eq}\left(\pmb{\sigma}_{n+1}\right):\pmb{\Xi}:\partial_{\pmb{\sigma}}\sigma_{eq}\left(\pmb{\sigma}_{n+1}\right) + \partial_{\Delta\gamma} Y\left(\bar\varepsilon^p_{n+1}\right)}\]

\[d\pmb{\sigma}_{n+1} = \left[\pmb{\Xi} - \frac{\partial_{\pmb{\sigma}}\sigma_{eq}\left(\pmb{\sigma}_{n+1}\right):\pmb{\Xi} \otimes \partial_{\pmb{\sigma}}\sigma_{eq}\left(\pmb{\sigma}_{n+1}\right):\pmb{\Xi} }{\partial_{\pmb{\sigma}}\sigma_{eq}\left(\pmb{\sigma}_{n+1}\right):\pmb{\Xi}:\partial_{\pmb{\sigma}}\sigma_{eq}\left(\pmb{\sigma}_{n+1}\right) + \partial_{\Delta\gamma} Y\left(\bar\varepsilon^p_{n+1}\right)}\right] :d\pmb{\varepsilon}_{n+1} \]



\end{document}
